\chapter[Atlas celular y reconstrucciones de masa]{Atlas celular y estatuto genealógico de las reconstrucciones de masa}
\label{chap:atlas-genealogia-masas}

El punto de partida es la jerarquía exacta del atlas: \(81\) celdas,
\(56\) familias de ruta y \(13\) multisecciones, con inversión central,
centro único y reparto residual \(12+12+12\). Para cada extensión de una
multisección, la evaluación sigue la cadena
\[
\text{extensión superviviente}\longrightarrow\text{ruta observable}
\longrightarrow\text{registro}\longrightarrow\text{firma}
\longrightarrow\chi_{\rm form}.
\]
La especialización global, el carácter positivo y las ocho torres actúan sólo
después; el sector nulo permanece separado de
\(\mathbb R_{>0}\). Sobre esta arquitectura se estudian, como objetos
subordinados y de tipo distinto, la restricción histórica calibrada de diez
filas y la decodificación inversa de doce pares.

El alcance de la arquitectura comprende la clasificación y reproducción de
las masas de las familias del Modelo Estándar y sus extensiones exóticas o
teóricas; se define por el atlas, las multisecciones y la cadena interna, con
independencia del número de filas de una comparación externa. Cada realización
concreta debe indicar, sin embargo, si la firma
procede de una ruta, de una reconstrucción calibrada o de una interfaz
prospectiva. Esta distinción preserva a la vez la amplitud del dominio y la
dirección causal de cada evaluación.

El conjunto de datos v11 se conserva, con su procedencia y huella, en el
material suplementario. Un verificador independiente reconstruye los censos y
retículos utilizados a continuación. Esta trazabilidad no convierte los
nombres físicos candidatos en teoremas de realización.

\section{Arquitectura interna y niveles de realización}

Conviene fijar de entrada las aplicaciones que aparecen en la cadena. Sean
\[
 \mathcal C_{81}=I_9\times I_9,
 \qquad I_9=\{1,\ldots,9\},
\]
el atlas celular, 
\(
 \mathcal Q_{10}=\{\mu,\pi^0,\pi^+,K^+,K^0,p,n,\tau,W,Z\}
\)
la restricción histórica calibrada de diez filas y
\(
 \mathcal Q_{12}^{(0)}=\{e,\mu,\tau,u,d,s,c,b,t,W,Z,H\}
\)
el conjunto de la tabla elemental de doce especies.

Los datos publicados definen, por separado, el mapa celular interno, la
restricción calibrada de pares y el carácter:
\[
 A_{\rm raw}:\mathcal C_{81}\longrightarrow\ZZ^5,
 \qquad
 (r,c)\longmapsto
 (n_A,s_C,k,\nu_{120},\nu_{270})_{r,c},
\]
\[
 \widehat\sigma_{10}:\mathcal Q_{10}
 \dashrightarrow\ZZ^5,
 \qquad (X,v_X)\ \text{fila calibrada},
\]
con conjunto de valores tabulados
\[
 \mathcal V_{10}:=\{v_X:(X,v_X)\ \text{es una fila calibrada}\},
\]
y el carácter positivo
\[
 \chi_{\rm mass}:\ZZ^5\longrightarrow\RR_{>0},
 \qquad v\longmapsto e^{\langle\lambda,v\rangle}.
\]
La notación discontinua no convierte las filas calibradas en una asignación
prospectiva desde la especie física: registra los pares de la tabla,
condicionados a objetivos de referencia. La tabla elemental de doce especies
define, en cambio, un decodificador inverso acotado
\[
 D_{\rm inv}:
 \{E_X^{\rm ref}:X\in\mathcal Q_{12}^{(0)}\}
 \longrightarrow\ZZ^2,
 \qquad
 E_X^{\rm ref}=m_X^{\rm ref}c^2
 \longmapsto(n_A,s_C)_X,
\]
obtenido minimizando un residuo respecto de una energía equivalente objetivo
dentro de una ventana finita.

La correspondencia prospectiva que falta parte de etiquetas y
representaciones físicas:
\[
 \mathcal Q_{\rm phys}
 \dashrightarrow
 \operatorname{MultiSec}_{\rm enr}(\Omega_{\rm H}^{\rm surv}).
 \tag{*a}
\]
Una vez dada una multisección enriquecida, la cadena interna y la carta
dimensional tienen la forma
\[
 \operatorname{MultiSec}_{\rm enr}(\Omega_{\rm H}^{\rm surv})
 \xrightarrow{\ \Gamma\ }
 \Gamma_{\TPK}^{[108]}
 \xrightarrow{\ L\ }
 \ZZ^5
 \xrightarrow{\ \chi_{\rm mass}\ }
 \RR_{>0}
 \xrightarrow{\ \varsigma_E\ }
 \mathcal U_E.
 \tag{*b}
\]
Las flechas hasta \(\RR_{>0}\) son internas, exactas o relativas una vez
fijada la extensión; \(\varsigma_E\) es una carta metrológica y no una flecha
interna adicional. El atlas construye la clasificación celular; el lector de extensiones
supervivientes produce la ruta enriquecida, el registro y la firma; el
carácter construye la evaluación adimensional. La flecha discontinua, rotulada
por su estatuto, denota la interfaz prospectiva que debe emparejar etiquetas y representaciones
físicas con multisecciones internas sin consultar la magnitud evaluada. No se
confunde con la tabla calibrada \(\widehat\sigma_{10}\) ni con el
decodificador inverso \(D_{\rm inv}\), cuyo dominio sí contiene energías
equivalentes de referencia \(E_X^{\rm ref}=m_X^{\rm ref}c^2\).

\begin{table}[htbp]
\centering
\small
\begin{tabularx}{\textwidth}{@{}
>{\raggedright\arraybackslash}p{0.24\textwidth}
>{\raggedright\arraybackslash}p{0.23\textwidth}
>{\raggedright\arraybackslash}X@{}}
\toprule
Objeto & Estatuto & Contenido certificado\\
\midrule
Atlas \(\mathcal C_{81}\) & exacto finito & celdas, familias, generaciones,
inversión, color residual y firmas crudas\\
Proyección física de capas & lectura candidata & nombres \(e,\mu,\tau,
\nu_e,\nu_\mu,\nu_\tau\) y sectores de quark sobre clases celulares\\
Restricción \(\widehat\sigma_{10}\) & \textsc{reconstrucción calibrada} & diez pares
etiqueta--vector condicionados a objetivos; no predice sin objetivo\\
Etapa \(E_4\) & \textsc{reconstrucción calibrada} & cuantización por entero más cercano
del residuo posterior a \(E_3\), dado el objetivo\\
Tabla \(\mathcal Q_{12}^{(0)}\) & decodificación inversa & representante de
Voronoi único en una ventana, elegido desde una energía equivalente de referencia\\
Carta de realización física & \textsc{interfaz física}/\textsc{programa abierto} & emparejamiento
prospectivo de etiquetas con multisecciones enriquecidas, separado del
decodificador inverso\\
\bottomrule
\end{tabularx}
\caption{Tipos distintos de evidencia en la genealogía de las masas.}
\label{tab:estatutos-atlas-masas}
\end{table}

\section{El atlas nonádico completo}

La involución central del atlas es
\[
 \rho:\mathcal C_{81}\longrightarrow\mathcal C_{81},
 \qquad \rho(r,c)=(10-r,10-c).
\]
Posee un único punto fijo, \((5,5)\), y cuarenta órbitas de dos elementos. En
consecuencia, \(\mathcal C_{81}/\langle\rho\rangle\) contiene \(41\) órbitas.
El censo celular se descompone como
\[
 81=45_{\rm leptones}+36_{\rm quarks}
\]
y, refinando por familia,
\[
 81=25_{\ell^-}+20_{\nu}+20_{d\text{-tipo}}+16_{u\text{-tipo}}.
\]

\begin{theorem}[Censo de familias y generaciones del atlas]
\label{thm:censo-atlas-81}
El atlas congelado posee exactamente la distribución
\[
\begin{array}{c|c|c}
\text{familia}&\text{generación celular}&\text{número de celdas}\\ \hline
\ell^-&G0,G2,G4&1,8,16\\
\nu&G1,G2,G3,G4&2,4,6,8\\
d\text{-tipo}&G1,G2,G3,G4&2,4,6,8\\
u\text{-tipo}&G1,G3&4,12.
\end{array}
\]
La involución \(\rho\) conserva sector, familia y generación. El centro
\((5,5)\) es la única celda \(G0\) y porta la etiqueta posicional
\(e_{\rm core}\).
\end{theorem}
\begin{proof}
La tabla contiene una fila para cada par de \(I_9^2\), sin duplicados. El
recuento por las columnas \texttt{sector}, \texttt{family} y
\texttt{generation} produce los valores indicados. Para cada fila
\((r,c)\), la fila \((10-r,10-c)\) posee los mismos campos de familia y
generación. La ecuación \((r,c)=\rho(r,c)\) fuerza \(r=c=5\); la fila central
es la única de generación cero.
\end{proof}

La palabra ``generación'' en este teorema nombra un estrato combinatorio del
atlas. Una identificación con las generaciones físicas exige un cociente de
sabor adicional. Esto resulta especialmente visible en el sector \(u\)-tipo:
el atlas v11 contiene únicamente \(G1\) y \(G3\), de modo que sus propias
etiquetas no proporcionan por sí solas una biyección con \(u,c,t\).
Asimismo, el punto fijo \((5,5)\) y su etiqueta interna \(e_{\rm core}\)
permanecen distinguidos dentro del atlas, pero no se identifican
automáticamente con el electrón físico: esa identificación pertenece a la
\textsc{interfaz física}/\textsc{programa abierto}.

\subsection{Color residual y conjugación central}

En el sector de quarks se define
\[
 \kappa(r,c)=r-c\pmod3,
\]
con la convención
\[
 0\longmapsto R,
 \qquad 1\longmapsto G,
 \qquad 2\longmapsto B.
\]
Como
\[
 \kappa(\rho(r,c))=(10-r)-(10-c)=-(r-c)\pmod3,
\]
la inversión central fija \(R\) e intercambia \(G\leftrightarrow B\).

\begin{proposition}[Balance de color residual]
\label{prop:color-atlas}
Las \(36\) celdas de quark se reparten de manera uniforme:
\[
 \#\kappa^{-1}(R)=\#\kappa^{-1}(G)=\#\kappa^{-1}(B)=12.
\]
La involución central actúa por conjugación del residuo de color.
\end{proposition}
\begin{proof}
El recuento de las celdas de quark en cada clase de \(r-c\bmod3\) da doce
elementos. La segunda afirmación es la identidad modular anterior.
\end{proof}

El resultado establece una clasificación residual ternaria exacta. El reparto
\(12+12+12\), por sí solo, no produce una representación. Fijada además la
carta qutrit de fase, el \cref{chap:algebra-color-qutrit} construye el par de
Weyl, \(\mathfrak{sl}_3(\mathbb C)\), su forma real compacta y una teoría gauge
finita. El entrelazador con el color físico y el límite cromodinámico continuo
permanecen como interfaz.

\subsection{Lectura física candidata de las capas}

El atlas propone la correspondencia de cociente siguiente:
\[
\begin{array}{c|c@{\qquad}c|c}
\text{sector celular}&\text{lectura candidata}&\text{celdas}&\text{estatuto}\\ \hline
\ell^-\,G0&e^-&1&\text{centro}\\
\ell^-\,G2&\mu^-&8&\text{capa}\\
\ell^-\,G4&\tau^-&16&\text{capa}\\
\nu\,G1&\nu_e&2&\text{capa neutra}\\
\nu\,G2&\nu_\mu&4&\text{capa neutra}\\
\nu\,G3&\nu_\tau&6&\text{capa neutra}\\
\nu\,G4&\nu_{\rm est/ex}&8&\text{capa estéril/exótica}\\
d\text{-tipo}\,G1,G2,G3,G4&\text{sectores de tipo abajo}&2,4,6,8&\text{requiere cociente}\\
u\text{-tipo}\,G1,G3&\text{sectores de tipo arriba}&4,12&\text{requiere cociente}.
\end{array}
\]
La tabla es una relación entre capas celulares y nombres físicos candidatos;
no es todavía una función que elija una ruta individual por especie. En
particular, las \(8\) celdas de la capa muónica y las \(16\) de la capa tauónica
requieren una orientación, una sección o un cociente adicional para producir
un único representante.

\section{Firmas crudas y retículo de masa}

Cada celda del archivo congelado porta una firma cruda
\[
 A_{\rm raw}(r,c)=(n_A,s_C,k,\nu_{120},\nu_{270})\in\ZZ^5.
\]
Las \(81\) celdas producen \(52\) vectores distintos. Sea
\[
 \Lambda_{\rm raw}
 =\left\langle A_{\rm raw}(r,c):(r,c)\in\mathcal C_{81}\right\rangle_\ZZ.
\]
Sea asimismo \(v_0=A_{\rm raw}(5,5)\) y
\[
 \Lambda_{\rm raw}^{0}
 =\left\langle A_{\rm raw}(r,c)-v_0\right\rangle_\ZZ.
\]

\begin{theorem}[Retículo crudo del atlas]
\label{thm:reticulo-crudo-atlas}
Ambos retículos tienen factores invariantes
\[
 \operatorname{SNF}(\Lambda_{\rm raw})
 =\operatorname{SNF}(\Lambda_{\rm raw}^{0})
 =(2,2,2,2,12).
\]
En particular, su índice en \(\ZZ^5\) es \(192\).
\end{theorem}
\begin{proof}
Las celdas \(22,23,39,56,57\) proporcionan una base de determinante absoluto
\(192\). La regla de Cramer verifica que cada una de las \(81\) firmas es una
combinación entera de esa base. Los divisores determinantales de la matriz
base son \(2,4,8,16,192\), que producen los factores
\((2,2,2,2,12)\). Para las diferencias respecto al centro, las celdas
\(13,14,37,54,74\) proporcionan otra base de determinante absoluto \(192\);
la misma comprobación de pertenencia y de divisores determinantales da el
segundo resultado.
\end{proof}

La restricción histórica calibrada contiene exactamente estas diez firmas:
\[
\begin{array}{c|rrrrr}
X&n_A&s_C&k&\nu_{120}&\nu_{270}\\ \hline
\mu&42&-3&8&0&2\\
\pi^0&44&-4&-25&1&-2\\
\pi^+&44&1&-2&2&-1\\
K^+&54&-1&17&-2&2\\
K^0&54&0&32&3&-2\\
p&59&0&9&1&0\\
n&59&1&-34&0&1\\
\tau&64&0&25&-1&6\\
W&94&-1&0&-2&4\\
Z&95&-1&-11&0&-4.
\end{array}
\]
Denotemos por \(\Lambda_{\rm mass}^{\Delta}\) el retículo generado por todas
sus diferencias.

\begin{theorem}[Separación de retículos y obstrucción directa]
\label{thm:obstruccion-atlas-masa}
Se cumplen
\[
 \operatorname{SNF}(\Lambda_{\rm mass}^{\Delta})=(1,1,1,1,2)
\]
y
\[
 A_{\rm raw}(\mathcal C_{81})
 \cap \mathcal V_{10}=\varnothing.
\]
Por tanto, las diez firmas de la restricción calibrada no son firmas de celdas crudas del
atlas v11. Además, el retículo de diferencias de masa contiene direcciones
unitarias que no genera el retículo crudo de índice \(192\).
\end{theorem}
\begin{proof}
El máximo común divisor de los menores de cada orden de la matriz de
diferencias de masa produce los factores \((1,1,1,1,2)\). La comparación
exhaustiva de los diez vectores con los \(52\) vectores crudos produce cero
coincidencias. La desigualdad de las formas normales de Smith impide
identificar los dos retículos como el mismo subretículo de \(\ZZ^5\).
\end{proof}

La diferencia de retículos demuestra que una completación de borde debe
añadir información si se desea pasar del atlas grueso a las firmas refinadas.
No construye por sí sola esa completación. En particular, llamar ``holonomía''
a la firma de masa expresa un programa geométrico preciso sólo después de
definir la conexión, las aristas admisibles y la independencia respecto de la
ruta elegida.

\subsection{Obstrucción de la reducción unifilar autónoma}

La dinámica unifilar del \cref{chap:masas} define sin consultar especies ni
masas una aplicación
\[
 F_{108}^{\rm uni}:I_9^2\times\{\pm1\}^2\longrightarrow\ZZ^5.
\]
Su dominio posee \(324\) condiciones iniciales y su imagen contiene \(53\)
vectores.

\begin{corollary}[No factorización por los mapas crudos vigentes]
\label{cor:no-factorizacion-atlas}
Se tiene
\[
 \operatorname{im}F_{108}^{\rm uni}
 \cap\mathcal V_{10}=\varnothing.
\]
No existe, por consiguiente, una interfaz prospectiva
\[
 S_{\rm uni}:\mathcal Q_{10}
 \dashrightarrow
 I_9^2\times\{\pm1\}^2
\]
tal que
\[
 F_{108}^{\rm uni}\bigl(S_{\rm uni}(X)\bigr)=v_X
 \qquad\text{para cada fila calibrada }(X,v_X).
\]
La misma conclusión vale si se sustituye \(F_{108}^{\rm uni}\) por la firma
cruda de las \(81\) celdas.
\end{corollary}
\begin{proof}
La enumeración autónoma de las \(324\) condiciones produce los \(53\)
vectores indicados y cero coincidencias con las diez firmas. Una
factorización exigiría que cada firma perteneciera a la imagen, lo que
contradice la intersección vacía. El segundo caso es el
\cref{thm:obstruccion-atlas-masa}.
\end{proof}

El corolario excluye dos factorizaciones concretas a través de reducciones que
han descartado información, pero no afecta al lector definido sobre el estado
TPK enriquecido. La carta directa debe conservar la extensión superviviente,
la orientación y el registro; por construcción permanece separada de toda
decodificación que tome una energía equivalente de referencia como entrada.

\section{Carácter de masa y jerarquía de refinamientos}

Los invariantes \(\Dfour,s_{120},\zCorr\) y \(s_{180}\) son globales: se
calculan una vez a partir de los canales \(q_\pm\) y se emplean en todas las
especies. Para
\[
 v=(n_A,s_C,k,\nu_{120},\nu_{270})\in\ZZ^5
\]
definimos
\[
 E_0(v)=n_AA_{\rm rad}+\frac{s_C}{6}\Cstarrad,
\]
\[
 E_1=E_0+k\Dfour,
 \qquad
 E_2=E_1+\nu_{120}s_{120},
 \qquad
 E_3=E_2+\nu_{270}\zCorr.
\]
El carácter
\[
 \chi_{\rm mass}(v)=e^{E_3(v)}
\]
es positivo y satisface
\[
 \chi_{\rm mass}(v+w)=\chi_{\rm mass}(v)\chi_{\rm mass}(w).
\]
Condicionada a la firma \(v_X\) y a una escala electrónica equivalente
\(E_e=m_ec^2\), la energía equivalente evaluada es
\[
 E_X^{(3)}=m_X^{(3)}c^2
 =E_e\chi_{\rm mass}(v_X).
\]

La restricción histórica calibrada de diez filas produce, con la escala electrónica interna,
los errores RMS relativos
\[
\begin{array}{c|c@{\qquad}c|c}
\text{etapa}&\operatorname{RMS}&\text{etapa}&\operatorname{RMS}\\ \hline
E_0&2.224237203948622\times10^{-3}&
E_1&2.762046114809764\times10^{-5}\\
E_2&4.777221368999134\times10^{-6}&
E_3&4.795312133659499\times10^{-7}\\
E_4&2.229279390173106\times10^{-8}&&
\end{array}
\]
Estos RMS son diagnósticos de calibración condicionados a
\(E_X^{\rm ref}\); no constituyen predicciones sin objetivo.

La cuarta corrección se define por
\[
 E_4(X)=E_3(X)+\nu_{180}(X)s_{180}.
\]
El carácter global de \(s_{180}\) está demostrado; el estatuto de
\(\nu_{180}(X)\) es diferente.

\begin{theorem}[Cuantización inversa del residuo \(E_4\)]
\label{thm:e4-inversa-atlas}
Para las diez filas de \(\mathcal Q_{10}\),
\[
 \nu_{180}(X)=
 \nint\!\left(
 \frac{\log(E_X^{\rm ref}/E_X^{(3)})}{s_{180}}
 \right),
\]
y el vector resultante, en el orden de la tabla, es
\[
 (2,7,-7,-1,-4,5,-4,-5,6,9).
\]
En consecuencia,
\[
 \left|
 \log\frac{E_X^{\rm ref}}{E_X^{(3)}}
 -\nu_{180}(X)s_{180}
 \right|\leq\frac{|s_{180}|}{2}.
\]
\end{theorem}
\begin{proof}
El entero se obtiene redondeando el cociente residual de cada energía
equivalente de referencia. La desigualdad es la propiedad del entero más cercano. La
evaluación de los diez cocientes produce el vector declarado y el RMS de
\(E_4\).
\end{proof}

El teorema demuestra que los residuos de \(E_3\) se alinean con una malla
global extraordinariamente fina. Como \(E_X^{\rm ref}\) aparece en la
definición de \(\nu_{180}\), esta etapa es reconstructiva. Para convertirla
en predicción prospectiva se requiere un selector de \(\nu_{180}\) obtenido
desde la especie o la ruta antes de leer la energía equivalente objetivo.

\section{Electrón: escala interna, origen de acción y paridad}

El invariante electrónico activo es
\[
 \mathfrak e_0=
 \frac{\sqrt3}{4}
 \left(e^{\varphi/\pi^2}-22\alphaHMT^3\right)
 (1+15\Dfour),
\]
con valor
\[
 \mathfrak e_0=
 0.510998922488066072509870877191349\ldots.
\]
Es una salida adimensional hasta elegir el transductor de energía. Bajo la
carta en MeV usada en la tabla se define la energía equivalente
\[
 E_e^{\rm HMT}=m_e^{\rm HMT}c^2
 =\mathfrak e_0\,\mathrm{MeV},
\]
y, con la referencia CODATA 2022 \cite{CODATA2022}, la comparación expresada
en MeV es
\[
 \frac{E_e^{\rm HMT}-E_e^{\rm CODATA}}{\mathrm{MeV}}
 =\mathfrak e_0-0.51099895069
 =-2.82019339\,10^{-8},
\]
con diferencia relativa aproximada
\(-5.519\,10^{-8}\). La energía equivalente electrónica externa
\(E_e^{\rm ext}=m_e^{\rm ext}c^2
=0.51099895\,\mathrm{MeV}\), utilizada en la tabla elemental histórica, es
una normalización distinta y no debe sustituir silenciosamente a
\(E_e^{\rm HMT}\).

Para estudiar el retículo se introduce además la firma de origen
\[
 v_e=(0,0,0,0,0).
\]
Esta convención de origen no afirma que la firma cruda de la celda \((5,5)\)
sea nula: de hecho, la celda central del atlas v11 porta una firma cruda
distinta. Se trata de dos niveles diferentes, posición celular y origen del
grafo de acción.

Sea
\[
 \ell(n_A,s_C,k,\nu_{120},\nu_{270})
 =s_C+\nu_{120}\pmod2.
\]

\begin{theorem}[Paridad no electrónica y extensión por el origen]
\label{thm:paridad-electron-atlas}
Las diez firmas de \(\mathcal V_{10}\) satisfacen \(\ell(v_X)=1\). Sus
diferencias generan
\[
 \ker\ell\subset\ZZ^5,
 \qquad
 \operatorname{SNF}=(1,1,1,1,2).
\]
Las diferencias desde el origen electrónico \(v_e\) generan \(\ZZ^5\), con
\[
 \operatorname{SNF}=(1,1,1,1,1).
\]
\end{theorem}
\begin{proof}
La evaluación de \(s_C+\nu_{120}\) da uno módulo dos para cada firma. Todas
las diferencias pertenecen, por ello, al núcleo. Sus divisores
determinantales dan índice dos; como \(\ell\) es sobreyectiva, el retículo de
diferencias coincide con \(\ker\ell\). Los diez vectores considerados como
aristas desde \(v_e=0\) poseen forma normal de Smith unitaria y generan el
retículo completo.
\end{proof}

Este teorema demuestra la función algebraica del origen electrónico una vez
dadas las firmas. No produce las firmas desde la celda central ni convierte
la etiqueta \(e_{\rm core}[5,5]\) en un selector de masa.

\section{Doce especies y teorema de decodificación inversa}

La tabla elemental histórica utiliza energías equivalentes:
\[
 E^{(0)}(n,s)=m^{(0)}(n,s)c^2
 =E_e^{\rm ext}
 \exp\!\left(nA_{\rm rad}+\frac{s}{6}\Cstarrad\right),
 \qquad E_e^{\rm ext}=m_e^{\rm ext}c^2
 =0.51099895\ \mathrm{MeV}.
\]
Para una energía equivalente de referencia \(E>0\), definimos la distancia logarítmica
\[
 d_E(n,s)=
 \left|
 \log\frac{E}{E_e^{\rm ext}}
 -nA_{\rm rad}-\frac{s}{6}\Cstarrad
 \right|
\]
en la ventana
\[
 \mathcal W=\{0,\ldots,120\}\times\{-12,\ldots,12\}.
\]

\begin{theorem}[Voronoi inverso acotado]
\label{thm:voronoi-doce-especies}
Para cada una de las doce energías equivalentes de referencia de la tabla, el
mínimo de \(d_E\) en \(\mathcal W\) es único. Los minimizadores son
\[
\begin{array}{c|rrrrrrrrrrrr}
X&e&\mu&\tau&u&d&s&c&b&t&W&Z&H\\ \hline
n_A&0&42&64&11&17&41&61&71&100&94&95&97\\
s_C&0&-3&0&7&8&-3&8&-5&-1&-1&-1&9.
\end{array}
\]
Las separaciones entre el mejor y el segundo mejor residuo son,
respectivamente,
\[
\begin{aligned}
 &(6.1777044,4.3915942,0.6382497,0.6995349,2.0918779,\
 4.8889835,\\
 &\qquad 3.8087408,1.3415669,5.1058939,6.0634351,
 3.7194917,1.6909808)\,10^{-3}.
\end{aligned}
\]
\end{theorem}
\begin{proof}
La ventana contiene \(121\cdot25=3025\) puntos. Para cada energía objetivo se evalúan y
ordenan exactamente esos \(3025\) residuos. El primer punto es el declarado
y su distancia es estrictamente menor que la del segundo; las diferencias
son las separaciones mostradas.
\end{proof}

El teorema prueba unicidad de decodificación dentro de la ventana y métrica
declaradas. Su dirección lógica es
\[
 E_X^{\rm ref}=m_X^{\rm ref}c^2\longmapsto(n_A,s_C)_X.
\]
No prueba una correspondencia prospectiva desde la etiqueta física \(X\)
hacia \((n_A,s_C)_X\) sin el dato de energía equivalente objetivo. Esta
distinción explica simultáneamente la estabilidad de
los pares publicados y su estatuto reconstructivo.

\section{Transducción de unidades}

Todos los exponentes y caracteres anteriores son adimensionales. Sea
\(\mathcal U_E\) una recta real de magnitudes de energía. Elegir una unidad
\(u_E\in\mathcal U_E\setminus\{0\}\) define una carta
\[
 \varsigma_{u_E}:\RR\longrightarrow\mathcal U_E,
 \qquad x\longmapsto x\,u_E.
\]

\begin{proposition}[Necesidad de una carta dimensional]
\label{prop:no-unidad-desde-adimensional}
Supóngase que los datos de entrada son invariantes bajo la acción de
reescalado \(u_E\mapsto\lambda u_E\), \(\lambda>0\). No existe una selección
natural de una magnitud de energía no nula a partir únicamente de esos datos.
\end{proposition}
\begin{proof}
Si \(v\in\mathcal U_E\) fuese una selección natural sin base elegida, debería
ser fija bajo todo reescalado: \(\lambda v=v\) para cada \(\lambda>0\). Tomando
\(\lambda\ne1\) se obtiene \(v=0\). Una salida no nula exige, por tanto, una
carta, una sección o un patrón dimensional adicional.
\end{proof}

La proposición no disminuye el contenido del carácter de masa: precisa dónde
entra la dimensión. La década de acción derivada por el determinante local
fija un orden decimal en la recta de acción, cuya carta puede expresarse en
\(\mathrm{J\,s}\). La identificación independiente de una base de energía con
MeV u otra unidad de energía corresponde a \(\varsigma_{u_E}\).

\section{Composición interna y carta de realización física}

El empalme matemático no termina en la firma cruda del atlas. El
\cref{chap:extension-ruta-caracter} construye la cadena
\[
\Omega_{\rm HMT}^{\rm surv}
\longrightarrow
\Gamma_{108}^{\rm enr}
\longrightarrow
\mathcal L_{108}
\xrightarrow{\ L_{\rm tick}\ }
\mathbb Z^5
\xrightarrow{\ \chi_{\rm mass}\ }
\mathbb R_{>0}.
\]
La extensión superviviente determina su sombra de ruta; la ruta enriquecida
conserva el registro; el funcional integral produce la firma de cinco
enteros; y el carácter evalúa esa firma sobre las torres globales. Ninguna de
estas flechas consulta una masa ni selecciona retrospectivamente una celda.

El atlas añade otra clasificación exacta. La aplicación interna
\[
 \mathcal R:\mathcal C_{81}\longrightarrow\mathcal R_{56}
\]
posee cincuenta y seis fibras de ruta, y el rango
\(\kappa\in\{0,1,2,3,4\}\) proporciona la elevación mínima que vuelve
inyectiva la pareja \((\mathcal R,\kappa)\). El cociente por familia y nivel
produce trece multisecciones. Doce de ellas son órbitas no puntuales de la
inversión central; la sección interna denominada electrónica es la única
excepción porque su fibra contiene el punto fijo \((5,5)\). El nombre interno
no realiza por sí solo el electrón físico.

La carta posterior tiene, por tanto, el tipo
\[
 \mathcal N_{\rm phys}^{\rm ind}
 \dashrightarrow
 \operatorname{MultiSec}_{\rm enr}
 \bigl(\Omega_{\rm HMT}^{\rm surv}\bigr).
\]
Su función pendiente es realizar etiquetas físicas individualizadas o
compuestas en las multisecciones enriquecidas y, una vez construidas, en
representaciones físicas compatibles. No reemplaza el
extractor de rutas ni el carácter de masa, y no exige que el nombre de una
especie elija un punto cuando la simetría interna sólo selecciona una órbita.

Quedan establecidos tres niveles. El atlas demuestra finitamente la jerarquía
\(81\to56\to13\). Fijada una extensión, las reglas declaradas demuestran el
empalme ruta \(\longrightarrow\) registro \(\longrightarrow\) firma
\(\longrightarrow\) evaluación; sólo la energía exige una carta metrológica.
Por separado, la restricción histórica calibrada de diez filas y la tabla
histórica de doce pares parten de objetivos de referencia y conservan,
respectivamente, los estatutos de reconstrucción calibrada y decodificación
inversa. La entrada prospectiva desde etiquetas físicas permanece como
\textsc{interfaz física}/\textsc{programa abierto}.
